\documentclass[10pt,a4paper]{amsart}
\usepackage[margin=19mm]{geometry}
\usepackage{amsmath,amssymb,mathtools}
\usepackage[scr=rsfs,cal=euler]{mathalfa}
\usepackage{xfrac}
\usepackage[unicode,hidelinks,bookmarks=false]{hyperref}
\newcommand{\ie}{i.e.\ }
\newcommand{\cf}{cf.\ }
\newcommand{\resp}{respectively\ }
\newcommand{\Subsection}{\subsection}
\providecommand{\nobreakdash}{\nobreak\hbox{-}}
\setlength{\emergencystretch}{2em}
\multlinegap0pt
\usepackage{bm}
\usepackage{bbm}
\usepackage{dsfont}
\usepackage{xspace}
\usepackage{cancel}
\usepackage{mathtools}
\usepackage[shortlabels]{enumitem}
\usepackage{amsthm}
\makeatletter
\@ifundefined{theorem}{\newtheorem{theorem}{Theorem}}{}
\@ifundefined{assumption}{\newtheorem{assumption}[theorem]{Assumption}}{}
\@ifundefined{lemma}{\newtheorem{lemma}[theorem]{Lemma}}{}
\@ifundefined{proposition}{\newtheorem{proposition}[theorem]{Proposition}}{}
\makeatother
\DeclareMathOperator*{\argmax}{arg\,max}
\title[Resource Extraction from a Stochastic Population with Harves]{Resource Extraction from a Stochastic Population with Harvesting Quotas}
\author{Erik Ekstr\"om, Dante Mata, Harto Saarinen}
\date{Source: arXiv:2608.17778v1 (CC BY 4.0)}
\begin{document}
\maketitle

\noindent\emph{Source and licence.} Resource Extraction from a Stochastic Population with Harvesting Quotas. Version arXiv:2608.17778v1 (CC BY 4.0), distributed under the Creative Commons Attribution 4.0 International licence, as recorded in the arXiv licence metadata for this exact version and shown on the abstract page https://arxiv.org/abs/2608.17778v1. Full rights notice: licenses/arxiv-2608.17778-CC-BY-4.0.txt.\par\smallskip

\noindent\textit{Editorial scope.} Counted unit: the Proposition whose source label is \texttt{Prop: k1k2ordering} in the article (source0.tex lines 688--696, source label \texttt{Prop: k1k2ordering}), delivered together with the article's own complete original proof of it (source0.tex lines 697--756, one \texttt{proof} environment of 60 lines). No mathematics is written, repaired or re-derived: statement and proof are the article's own text line for line. Required context retained uncounted: Ass:MainAssumptions (lines 153--171); hit 0 (lines 395--405); Lemma: Hshape (lines 484--486). Every definition, display and label of those lines is retained unchanged. Omitted from this package and not redistributed: 1--152, 172--394, 406--483, 487--687, 757--881. Nothing inside a retained excerpt is omitted or altered: every retained body is delivered exactly as the article's lines are written, so no omission receipt of any kind accompanies this package. Citations: the retained excerpts carry no citation command at all, so no bibliography entry of the article is transcribed and no reference is redistributed. Reference adaptation: none was needed and none was made; every reference inside the delivered unit resolves inside it, so no unresolved reference or citation is produced. Printed slips kept literal: no printed word, symbol, hypothesis or number of the counted statement or of its proof was altered. Layout and class: the article's own class is \texttt{article} with packages including amsmath, amsthm, bbm, amssymb, amsfonts, graphicx, datetime, enumerate, float, natbib, hyperref, geometry, xcolor; the wrapper is the lane's pinned \texttt{amsart} editorial wrapper at 10pt on A4, which restates the theorem-like environments the retained text needs and adds the article's own macro definitions that the retained text needs (\texttt{\textbackslash argmax}), with extra wrapper packages (bm, bbm, dsfont, xspace, cancel, mathtools, enumitem). The delivered numbering is the unit's own. Assets: no retained span contains a figure environment, a graphics-inclusion command, a PDF-page inclusion, a table or a TikZ picture; no image file is copied into this package, the package declares no asset dependency and no image file is redistributed with it. Licence: version arXiv:2608.17778v1 of this article is distributed under the Creative Commons Attribution 4.0 International licence, as recorded in the arXiv licence metadata for this exact version and shown on the abstract page https://arxiv.org/abs/2608.17778v1. A full rights notice is delivered at licenses/arxiv-2608.17778-CC-BY-4.0.txt. Characters: every retained excerpt is pure ASCII, so the delivered engine, XeLaTeX with its default Latin Modern text font, reports no missing character.
\begin{assumption} \label{Ass:MainAssumptions} 
The coefficients $\mu$ and $\sigma$ are locally Lipschitz continuous with $\sigma^2(x)>0$ for $x>0$, and such that 
    \begin{enumerate}     
            \item[(i)]
            the boundary point $\infty$ is natural, and 0 is natural or entrance-not-exit;
            \item[(ii)]
        there exists an $x^* =\argmax_x \{ \mu(x) \}>0$ so that $\mu(x)$ is increasing on $[0, x^*)$ and decreasing on $(x^*, \infty)$, with
        \[ \mu(0) \geq 0 \quad\&\quad\lim_{x \to \infty} \mu(x) < 0;\]
        \item[(iii)]
        $\int_0^{\infty} m(dy) <\infty$;
        \item[(iv)]
        $S'(\infty):=\lim_{x\to\infty}S'(x)=\infty$.
    \end{enumerate}
    Moreover, 
      \begin{enumerate}     
            \item[(v)]
            $k:[0,\infty)\to[0,\infty)$ is non-decreasing, and $k(x)>0$ for $x>0$.
              \end{enumerate}
\end{assumption}

\begin{assumption}\label{hit 0}
We have that
\begin{enumerate}[(i)]
    \item 
    $\int_0^\infty m_0(dy)<\infty$;
    \item 
    the quantity 
\[\overline\mu:=\frac{\int_0^{\infty}\mu(y)m_0(dy)}{\int_0^\infty m_0(dy)}\]
    satisfies $\mu(0)<\overline\mu$.
\end{enumerate}
\end{assumption}

\begin{lemma} \label{Lemma: Hshape}
    There exists a point $x_H \in(0,b^*)$ such that the function $H$ is increasing and satisfies $H> \mu$ on $(0,x_H)$, and decreasing with $H<\mu$ on $(x_H,\infty)$. Consequently, $b^* \in (x_H, x')$.
\end{lemma}

\begin{proposition} \label{Prop: k1k2ordering}
    Let $k_1,k_2: (0,\infty) \to (0,\infty)$ be two extraction bounds satisfying Assumptions~\ref{Ass:MainAssumptions} and \ref{hit 0}, with $k_1(x) \leq k_2(x)$ for all $x \geq 0.$
    
    For $i=1,2$, let
    \begin{equation*}
        H_i(x) := \frac{\int_x^\infty \mu(y)m_{i}(dy)}{\int_x^\infty m_{i}(dy)},
    \end{equation*}
    where $m_{i}$ is the speed measure corresponding to the drift $\mu-k_i$, and let $b_i$ denote the unique positive solution of $G(b_i) = H_i(b_i).$ Then $b_1 \leq b_2.$
\end{proposition}

\begin{proof}
    Let 
    \begin{equation*}
       \lambda_i(x) := \frac{m_{i}'(x)}{\int_x^\infty m_{i}(dy)}\quad\& \quad  \Delta(x) := H_2(x)-H_1(x).
    \end{equation*}
By the definition of the speed measure we have $m_{2}'(y) = q(y)m_{1}'(y),$ where
    \begin{equation*}
        q(y) = \exp{\left(-\int^{y}_{c_0}\frac{2(k_2(z)-k_1(z))}{\sigma^2(z)}dz \right)}.
    \end{equation*}
    As $k_2(z)-k_1(z) \geq 0$ for all $z>0$, the function $q$ is nonincreasing. Therefore we find
    \begin{equation}\label{lambdaineq}
        \lambda_2(x) = \frac{q(x)m_{1}'(x)}{\int_x^\infty q(y)m_{1}'(y)dy} \geq \frac{m_{1}'(x)}{\int_x^\infty m_{1}'(y)dy} = \lambda_1(x).
    \end{equation}
    
    Next, differentiation yields $H_i'(x) =  \lambda_i(x) \big(H_i(x)-\mu(x)\big)$, and thus
    \begin{equation*}
        \Delta'(x)-\lambda_2(x)\Delta(x) = \big( \lambda_2(x)-\lambda_1(x)\big)\big(H_1(x)-\mu(x) \big).
    \end{equation*}
    Let $x_H^1$ be the point at which $H_1-\mu$ changes sign, cf. Lemma~\ref{Lemma: Hshape}. Then
    \begin{equation*}
        H_1(x)-\mu(x) \leq 0 \quad \text{ for all } x \geq x_H^1.
    \end{equation*}
    Combining this with \eqref{lambdaineq}, we obtain for all $x \geq x_H^1$ that 
    \begin{equation} \label{Eq: DeltaIneq}
        \Delta'(x)-\lambda_2(x)\Delta(x) \leq 0.
    \end{equation}
    Consequently,
    \begin{equation*}
        \frac{d}{dx} \left( \Delta(x)\int_x^\infty m_{2}(dy) \right) = \int_x^\infty m_{2}(dy) \big(\Delta'(x)-\lambda_2(x)\Delta(x)\big) \leq 0,
    \end{equation*}
    and thus the function $x \mapsto \Delta(x) \int_x^\infty m_{2}(dy)$ is nonincreasing on $[x_H^1,\infty)$. 
    
    Now, since $q$ is nonincreasing, the ratio 
    \begin{equation*}
        \frac{\int_x^\infty m_{2}(dy)}{\int_x^\infty m_{1}(dy)}
        =\frac{\int_x^\infty q(y)m_{1}(dy)}{\int_x^\infty m_{1}(dy)}
    \end{equation*}
    is bounded by $q(x_H^1)$ on $[x_H^1, \infty)$. Noting that 
    \[\int_x^\infty \mu(y)m_{i}(dy) \to 0\] 
    as $x \to \infty$, it follows that
    \begin{equation*}
        \Delta(x)\int_x^\infty m_{2}(dy) = 
        \int_x^\infty \mu(y)m_{2}(dy)-\int_x^\infty \mu(y)m_{1}(dy)(x)\frac{\int_x^\infty m_{2}(dy)}{\int_x^\infty m_{1}(dy)} \to 0
    \end{equation*}
    as $x \to \infty$. Hence, we have 
    \begin{equation*}
        \Delta(x) \int_x^\infty m_{2}(dy) \geq 0 \quad \text{ for all } x \geq x_H^1,
    \end{equation*}
    which implies
    \begin{equation*}
        \Delta(x) \geq 0 \quad \text{ for all } x \geq x_H^1.
    \end{equation*}
    Consequently, $H_2(x) \geq H_1(x)$ for all $ x \geq x_H^1$.
    
    Finally, since the unique root $b_1 \in (x_H^1,x')$, we have 
    \begin{equation*}
        H_2(b_1) \geq H_1(b_1) = G(b_1).
    \end{equation*}
    As the equation $G(x) = H_2(x)$ has the unique solution $b_2$, and since $G(x)-H_2(x)$ changes sign from negative to positive at that point, it follows that $b_1 \le b_2.$
\end{proof}

\end{document}
